\problemstart
\addcontentsline{toc}{section}{1.3　極座標における速度と加速度}

\begin{problemBox}{問題 1.3\quad 極座標における速度と加速度\hspace{1.2em}{\small\mdseries 〔標準〕}}
平面上を運動する質点の位置を記述するため，二次元デカルト座標の単位ベクトルを $\boldsymbol{i}, \boldsymbol{j}$ とし，平面極座標の動径方向の単位ベクトルを $\boldsymbol{e}_r$ ，方位角方向の単位ベクトルを $\boldsymbol{e}_\theta$ とする。位置ベクトルが $\boldsymbol{r} = r \boldsymbol{e}_r$ で表されるとき，以下の問いに答えよ。

\begin{enumerate}
  \item デカルト座標の単位ベクトルを用いて極座標の単位ベクトルは次のように表される。
  \begin{equation*}
  \boldsymbol{e}_r = \boldsymbol{i} \cos\theta + \boldsymbol{j} \sin\theta
  \end{equation*}
  \begin{equation*}
  \boldsymbol{e}_\theta = -\boldsymbol{i} \sin\theta + \boldsymbol{j} \cos\theta
  \end{equation*}
  これらの単位ベクトルの時間微分 $\frac{d\boldsymbol{e}_r}{dt}$ および $\frac{d\boldsymbol{e}_\theta}{dt}$ を計算し， $\boldsymbol{e}_r, \boldsymbol{e}_\theta$ および角速度 $\dot{\theta} = \frac{d\theta}{dt}$ を用いて表せ。

  \item 位置ベクトル $\boldsymbol{r} = r \boldsymbol{e}_r$ を時間 $t$ で微分することにより，速度ベクトル $\boldsymbol{v}$ を求めよ。また，動径方向成分 $v_r$ と方位角方向成分 $v_\theta$ をそれぞれ示せ。

  \item 速度ベクトル $\boldsymbol{v}$ をさらに時間 $t$ で微分することにより，加速度ベクトル $\boldsymbol{a}$ を求めよ。また，動径方向成分 $a_r$ と方位角方向成分 $a_\theta$ をそれぞれ示せ。
\end{enumerate}
\end{problemBox}

\solutionheading
\begin{solutionparts}

\solutionpart{(1)}
時間変化しない固定された単位ベクトル $\boldsymbol{i}, \boldsymbol{j}$ に対し，極座標の単位ベクトルは角度 $\theta$ のみを介して時間に依存します。合成関数の微分公式を適用して時間微分を計算します。
\begin{align*}
\frac{d\boldsymbol{e}_r}{dt}
  &= \frac{d\boldsymbol{e}_r}{d\theta}\frac{d\theta}{dt} \\
  &= (-\boldsymbol{i}\sin\theta+\boldsymbol{j}\cos\theta)\dot{\theta}
\end{align*}
括弧の内部は方位角方向の単位ベクトル $\boldsymbol{e}_\theta$ そのものであるため，次のようになります。
\begin{equation*}
\frac{d\boldsymbol{e}_r}{dt} = \dot{\theta} \boldsymbol{e}_\theta
\end{equation*}
同様に，もう一方の単位ベクトルについても計算します。
\begin{align*}
\frac{d\boldsymbol{e}_\theta}{dt}
  &= \frac{d\boldsymbol{e}_\theta}{d\theta}\frac{d\theta}{dt} \\
  &= (-\boldsymbol{i}\cos\theta-\boldsymbol{j}\sin\theta)\dot{\theta} \\
  &= -(\boldsymbol{i}\cos\theta+\boldsymbol{j}\sin\theta)\dot{\theta}
\end{align*}
括弧の内部は動径方向の単位ベクトル $\boldsymbol{e}_r$ であるため，次のようになります。
\begin{equation*}
\frac{d\boldsymbol{e}_\theta}{dt} = -\dot{\theta} \boldsymbol{e}_r
\end{equation*}

\solutionpart{(2)}
位置ベクトル $\boldsymbol{r} = r \boldsymbol{e}_r$ の両辺を時間 $t$ で微分します。積の微分公式を適用すると次のようになります。
\begin{equation*}
\boldsymbol{v} = \frac{d\boldsymbol{r}}{dt} = \frac{dr}{dt} \boldsymbol{e}_r + r \frac{d\boldsymbol{e}_r}{dt}
\end{equation*}
ここに (1) の結果である $\frac{d\boldsymbol{e}_r}{dt} = \dot{\theta} \boldsymbol{e}_\theta$ を代入し，時間微分をドット記号で表現します。
\begin{equation*}
\boldsymbol{v} = \dot{r} \boldsymbol{e}_r + r \dot{\theta} \boldsymbol{e}_\theta
\end{equation*}
各単位ベクトルの係数がそれぞれの成分に対応するため，求める成分は次のようになります。
\begin{equation*}
v_r = \dot{r}
\end{equation*}
\begin{equation*}
v_\theta = r \dot{\theta}
\end{equation*}

\solutionpart{(3)}
(2) で得られた速度ベクトル $\boldsymbol{v} = \dot{r} \boldsymbol{e}_r + r \dot{\theta} \boldsymbol{e}_\theta$ を時間 $t$ でさらに微分します。和の各項に対して積の微分公式を適用します。
\begin{align*}
\boldsymbol{a}
  &= \frac{d\boldsymbol{v}}{dt} \\
  &= \left(\ddot{r}\boldsymbol{e}_r
     +\dot{r}\frac{d\boldsymbol{e}_r}{dt}\right) \\
  &\quad +\left(\dot{r}\dot{\theta}\boldsymbol{e}_\theta
     +r\ddot{\theta}\boldsymbol{e}_\theta
     +r\dot{\theta}\frac{d\boldsymbol{e}_\theta}{dt}\right)
\end{align*}
この式に (1) で求めた単位ベクトルの微分式をそれぞれ代入します。
\begin{align*}
\boldsymbol{a}
  &= \ddot{r}\boldsymbol{e}_r
   + \dot{r}(\dot{\theta}\boldsymbol{e}_\theta)
   + \dot{r}\dot{\theta}\boldsymbol{e}_\theta \\
  &\quad +r\ddot{\theta}\boldsymbol{e}_\theta
   +r\dot{\theta}(-\dot{\theta}\boldsymbol{e}_r)
\end{align*}
動径方向の単位ベクトル $\boldsymbol{e}_r$ を含む項と，方位角方向の単位ベクトル $\boldsymbol{e}_\theta$ を含む項に整理してまとめます。
\begin{align*}
\boldsymbol{a}
  &= (\ddot{r}-r\dot{\theta}^2)\boldsymbol{e}_r \\
  &\quad +(\dot{r}\dot{\theta}+\dot{r}\dot{\theta}
     +r\ddot{\theta})\boldsymbol{e}_\theta
\end{align*}
等しい項を足し合わせて式を簡潔にします。
\begin{equation*}
\boldsymbol{a} = (\ddot{r} - r \dot{\theta}^2) \boldsymbol{e}_r + (r \ddot{\theta} + 2 \dot{r} \dot{\theta}) \boldsymbol{e}_\theta
\end{equation*}
各単位ベクトルの係数を取り出すことで，加速度の各成分が得られます。
\begin{equation*}
a_r = \ddot{r} - r \dot{\theta}^2
\end{equation*}
\begin{equation*}
a_\theta = r \ddot{\theta} + 2 \dot{r} \dot{\theta}
\end{equation*}
動径方向の第2項は遠心力に関連する向心加速度の項であり，方位角方向の第2項はコリオリの力に関連する加速度の項に相当します。

\end{solutionparts}
